\subsection{Continuous positive complex Hessians}
\label{subsec:continuous-hessian-adjoints}

The smooth recovery theorem extends to a bounded solution whose complex
Hessian is continuous and everywhere positive. In this case one can
vary the potential in both directions, and the resulting densities
force every supporting probability to annihilate the same elliptic
operator. The Sobolev range theorem then gives uniqueness.

\Needspace{14\baselineskip}
\begin{theorem}\label{thm:continuous-hessian-adjoint-uniqueness}\index{Adjoint measure!continuous positive Hessian}
Let $n\geq2$, let $\theta=\omega$ be a Hermitian form, and let $\varphi=0$.
Suppose $F\in C(X)$ and a bounded solution $u\in\PSH(X,\omega)$ of
\eqref{eq:relative-normalization} has distributional complex Hessian
represented by a continuous form
\[
 \sigma=\omega+\ddc u>0
\]
everywhere on $X$. Then $\mathcal A_0(F)$ consists of one probability
$\nu_F$. It satisfies
\begin{align*}
 \int_X L_\sigma h\,\mathrm d\nu_F&=0\quad(h\in C^\infty(X)),\\
 DB_0(F)[H]&=-\int_XH\,\mathrm d\nu_F\quad(H\in C(X)).
\end{align*}
The derivative is taken in the G\^ateaux sense.
\end{theorem}
\begin{proof}
The continuous trace of $\ddc u$ is the Euclidean Laplacian of $u$, up
to the fixed normalization. Local constant-coefficient elliptic
regularity gives $u\in W^{2,p}$ for every finite $p$. Indeed, subtract
the Newtonian potential of a cutoff of this trace. Its second
high-frequency derivatives are order-zero multipliers, and the
low-frequency kernels are locally bounded; the remaining distribution
is harmonic. Sobolev embedding gives a $C^{1,\alpha}$ representative
for every $\alpha<1$, which agrees with the quasi-psh representative.

On a coordinate ball choose smooth $\rho$ with $\ddc\rho\geq\omega$.
The radial regularizations of the bounded psh function $u+\rho$ decrease
to it, and their complex Hessians converge uniformly on smaller balls.
Continuity of bounded products and the finite smooth-background
expansion identify $(\omega+\ddc u)^n$ with $\sigma^n$. For smooth real
$h$, the same argument applies to $u+th$, which is admissible for all
sufficiently small signed $t$ by uniform positivity. Define
\[
 F_t=F+\log\frac{(\sigma+t\ddc h)^n}{\sigma^n}.
\]
Its equation has the same normalization scalar as the equation for $u$,
so $B_0(F_t)=B_0(F)$. The supporting inequality gives
$\int_X(F_t-F)\,\mathrm d\nu\leq0$ for each
$\nu\in\mathcal A_0(F)$. Since
$F_t-F=tL_\sigma h+O(t^2)$ uniformly, using both signs yields stationarity.

Choose $p=4n+1$ and $q=p/(p-1)=1+1/(4n)$. By
\cref{cor:gaud-adjoint-lq}, every supporting probability has an $L^q$
density, since $q<n/(n-1)$. Smooth density in $W^{2,p}$ extends its
stationarity to every $W^{2,p}$ test. By
\cref{lem:continuous-hermitian-fredholm}, all these densities lie in a
one-dimensional space. Their integrals equal one, so they coincide.
Nonemptiness is \cref{thm:weak-adjoint-probabilities}, and the derivative
formula follows from \cref{prop:normalization-directional}.
\end{proof}

\subsection{Products with a Hermitian factor}

An explicit product construction exhibits the criterion for densities
with an arbitrary continuous modulus. Let $(Y,\kappa)$ be a connected
compact Hermitian $k$-fold, $k\geq1$, with $\int_Y\kappa^k=1$.
Denote by $\nu_\kappa$ the smooth stationary probability of
$L_\kappa=\operatorname{tr}_\kappa\ddc$, given by the Gauduchon multiplier
when $k\geq2$; for $k=1$ it is $\kappa$. Write $w$ for the
Fubini--Study area of mass one on $\mathbb P^1$. All factor forms and
functions below are pulled back to $X=\mathbb P^1\times Y$.

\begin{proposition}\label{prop:continuous-product-adjoint}\index{Adjoint measure!Hermitian product}
Let $a>0$, let $b$ be smooth and positive on $\mathbb P^1$, and let $A$
be continuous and positive there, with $\int_{\mathbb P^1}Aw=a$.
For
\[
 \omega=aw+b\kappa,\qquad \mathrm dV=w\wedge\kappa^k,
 \qquad F=\log((k+1)Ab^k),
\]
one has
\[
 c_0(F)=1,\qquad
 \mathcal A_0(F)=\left\{\frac Aa\,w\otimes\nu_\kappa\right\}.
\]
\end{proposition}
\begin{proof}
The Green solution of $\ddc u=(A-a)w$ on $\mathbb P^1$ has continuous
complex Hessian and is bounded. Its pullback yields
$\sigma=Aw+b\kappa>0$, and
$\sigma^{k+1}=(k+1)Ab^kw\wedge\kappa^k$.
Thus scalar uniqueness gives $c_0(F)=1$, and
\cref{thm:continuous-hessian-adjoint-uniqueness} gives a unique support.
For every smooth $h$ depending on all variables, the block decomposition
\[
 L_\sigma h=A^{-1}L_wh+b^{-1}L_\kappa h
\]
shows that $(A/a)w\otimes\nu_\kappa$ is stationary: the first term
integrates to zero by Stokes on $\mathbb P^1$, and the second by
stationarity on $Y$ and Fubini. This candidate has a bounded density
with respect to $\mathrm dV$. The one-dimensional annihilator of
\cref{lem:continuous-hermitian-fredholm} and probability normalization
therefore identify it with the supporting probability.
\end{proof}

\begin{example}\label{ex:non-dini-hopf-product-adjoint}\index{Dini continuity!failure with unique adjoint}\index{Hopf manifold!adjoint product example}
There is a non-K\"ahler threefold with a nonclosed Hermitian background
and a positive continuous density whose logarithm is neither Dini
continuous nor H\"older continuous of any positive order, but whose
supporting probability is explicit and unique.
\end{example}
\begin{proof}
Take the Hopf surface
$Y=(\mathbb C^2\setminus\{0\})/\langle z\mapsto2z\rangle$ and its
unit-volume metric $\kappa=c|z|^{-2}\ddc|z|^2$, where $c>0$ is chosen
so that $\int_Y\kappa^2=1$. Put $s=\log|z|^2$, $\alpha=\ddc s$ and
$\beta=\mathrm ds\wedge\dc s$, so that $\kappa=c(\alpha+\beta)$.
The differentiation in \eqref{eq:hopf-torsion} applies in this dimension.
Here $\alpha$ has rank one, so $\ddc\kappa=-c\alpha^2=0$ and
$\nu_\kappa=\kappa^2$. Polar coordinates identify $Y$ with
$S^3\times S^1$, hence $H^2(Y,\mathbb R)=0$. A closed positive form
would be exact and have positive square integral, contradicting Stokes.
Thus $Y$ is non-K\"ahler, and so is $\mathbb P^1\times Y$, by restriction.

Let $\tau\in[-1,1]$ be the height on $\mathbb P^1$, set $b=2+\tau$,
and put
\[
 f(t)=\begin{cases}
       1/\log(e/|t|),&0<|t|\leq1,\\
       0,&t=0,
      \end{cases}
\]
\[
 A=a+\varepsilon\left(f(\tau)-\int_{\mathbb P^1}f(\tau)w\right),
 \quad0<\varepsilon<a.
\]
Since $0\leq f\leq1$, this is a positive continuous function of mean
$a$. Along a normal segment to the equator, $|\tau|$ is comparable
to the distance $r$. The oscillation of $A$ is therefore bounded below
by a positive multiple of $1/\log(e/r)$. Its Dini integral
$\int_0^{r_0}\mathrm dr/(r\log(e/r))$ diverges, and this modulus is
larger than every $r^\gamma$, $\gamma>0$, as $r\downarrow0$.
The logarithm of $A$ has the same lower bound, because $A$ is uniformly
positive and bounded. The smooth term $2\log b$ contributes only $O(r)$,
so $F=\log(3Ab^2)$ has neither regularity. The mixed component
$\mathrm db\wedge\kappa$ shows that $\omega=aw+b\kappa$ is nonclosed.
By \cref{prop:continuous-product-adjoint}, its supporting probability
is $(A/a)w\wedge\kappa^2$.
\end{proof}

The theorem assumes a continuous positive complex Hessian for a bounded
solution. It does not establish that hypothesis for every continuous $F$;
the uniqueness question for arbitrary continuous data remains open.
